\documentclass[12pt]{article}
\usepackage[margin=1in]{geometry}
\usepackage{amsmath,amssymb}
\usepackage{enumitem}
\usepackage{array}
\usepackage{booktabs}

\begin{document}

\begin{center}
    {\Large \textbf{Chapter 7 and Chapter 8 Solutions}}
\end{center}

\section*{Chapter 7}

\begin{enumerate}[label=\textbf{\arabic*.}, leftmargin=*]

\setcounter{enumi}{1}
\item
\begin{enumerate}[label=\alph*.]
    \item The sample mean is the center of the interval, so
    \[
    \bar{x}=\frac{114.4+115.6}{2}=115.
    \]
    \item The interval $(114.4,115.6)$ has the $90\%$ confidence level. The higher confidence level will produce a wider interval.
\end{enumerate}

\setcounter{enumi}{3}
\item
\begin{enumerate}[label=\alph*.]
    \item
    \[
    58.3\pm \frac{1.96(3)}{\sqrt{25}}=58.3\pm1.18=(57.1,59.5).
    \]
    \item
    \[
    58.3\pm \frac{1.96(3)}{\sqrt{100}}=58.3\pm0.59=(57.7,58.9).
    \]
    \item
    \[
    58.3\pm \frac{2.58(3)}{\sqrt{100}}=58.3\pm0.77=(57.5,59.1).
    \]
    \item For an $82\%$ confidence interval,
    \[
    1-\alpha=.82\Rightarrow \alpha=.18\Rightarrow \alpha/2=.09,
    \]
    and $z_{.09}=1.34$. Therefore,
    \[
    58.3\pm \frac{1.34(3)}{\sqrt{100}}=(57.9,58.7).
    \]
    \item
    \[
    n=\left[\frac{2(2.58)3}{1}\right]^2=239.62\approx240.
    \]
\end{enumerate}

\setcounter{enumi}{9}
\item
\begin{enumerate}[label=\alph*.]
    \item When $n=15$, $2\lambda\sum X_i$ has a chi-squared distribution with $30$ df. From the $30$ df row of Table A.6, the critical values that capture lower and upper tail areas of $.025$ are $16.791$ and $46.979$. Thus
    \[
    \left(\frac{2\sum x_i}{46.979},\frac{2\sum x_i}{16.791}\right)
    \]
    is a $95\%$ confidence interval for
    \[
    \mu=\frac{1}{\lambda}.
    \]
    Since $\sum x_i=63.2$, the interval is
    \[
    (2.69,7.53).
    \]

    \item A $99\%$ confidence level requires critical values that capture area $.005$ in each tail of the chi-squared curve with $30$ df. These are $13.787$ and $53.672$, which replace $16.791$ and $46.979$ in part (a).

    \item
    \[
    V(X)=\frac{1}{\lambda^2}
    \]
    when $X$ has an exponential distribution, so the standard deviation is
    \[
    \frac{1}{\lambda},
    \]
    the same as the mean. Thus the interval in part (a) is also a $95\%$ confidence interval for the standard deviation of the lifetime distribution.
\end{enumerate}

\setcounter{enumi}{15}
\item
\begin{enumerate}[label=\alph*.]
    \item The boxplot shows a high concentration in the middle half of the data, with a narrow box width. There is a single outlier at the upper end, but that value is closer to the median, $55$ kV, than to the smallest sample observation.

    \item From software,
    \[
    \bar{x}=54.7,\qquad s=5.23.
    \]
    The $95\%$ confidence interval is
    \[
    \bar{x}\pm1.96\frac{s}{\sqrt{n}}
    =54.7\pm1.96\frac{5.23}{\sqrt{48}}
    =54.7\pm1.5=(53.2,56.2).
    \]
    We are $95\%$ confident that the true mean breakdown voltage under these conditions is between $53.2$ kV and $56.2$ kV. The interval is reasonably narrow, indicating that $\mu$ has been estimated fairly precisely.

    \item A conservative estimate of the standard deviation is
    \[
    \frac{70-40}{4}=7.5.
    \]
    To achieve a margin of error of at most $1$ kV with $95\%$ confidence,
    \[
    1.96\frac{s}{\sqrt{n}}\le1
    \Rightarrow
    n\ge\left[\frac{1.96s}{1}\right]^2
    =\left[\frac{1.96(7.5)}{1}\right]^2=216.09.
    \]
    Therefore, a sample of size at least $n=217$ is required.
\end{enumerate}

\setcounter{enumi}{19}
\item
\begin{enumerate}[label=\alph*.]
    \item With the numbers provided, a $99\%$ confidence interval for the parameter $p$ is
    \[
    \frac{.53+\frac{(2.58)^2}{2(2343)}\pm 2.58\sqrt{\frac{(.53)(.47)}{2343}+\frac{(2.58)^2}{4(2343)^2}}}{1+\frac{(2.58)^2}{2343}}
    =(.503,.556).
    \]
    We are $99\%$ confident that the proportion of all adult Americans who watched streamed programming was between $.503$ and $.556$.

    \item
    \[
    n=\left[\frac{2z_{\alpha/2}\sqrt{\hat{p}\hat{q}}}{w}\right]^2
    =\left[\frac{2(2.58)(.5)(.5)}{.05}\right]^2=665.64,
    \]
    so $n$ must be at least $666$.
\end{enumerate}

\setcounter{enumi}{25}
\item
Using the quadratic equation to solve the limits
\[
\frac{\bar{x}-\mu}{\sqrt{\mu/n}}=\pm z
\]
gives the solutions
\[
\mu=\left(\bar{x}+\frac{z^2}{2n}\right)\pm z\frac{\sqrt{z^2+4n\bar{x}}}{2n}.
\]
For the data provided, $n=50$ and
\[
\bar{x}=\frac{203}{50}=4.06.
\]
Using $z=1.96$, we are $95\%$ confident that the true value of $\mu$ is in the interval
\[
\left(4.06+\frac{1.96^2}{2(50)}\right)
\pm1.96\frac{\sqrt{1.96^2+4(50)(4.06)}}{2(50)}
=4.10\pm.56=(3.54,4.66).
\]
For large $n$, this is approximately equivalent to
\[
\bar{x}\pm z\frac{\sqrt{\bar{x}}}{\sqrt{n}}=\hat{\mu}\pm z\frac{\hat{\sigma}}{\sqrt{n}},
\]
since $\mu=\sigma^2$ for the Poisson distribution.

\setcounter{enumi}{29}
\item
\begin{enumerate}[label=\alph*.]
    \item $t_{.025,10}=2.228$
    \item $t_{.025,15}=2.131$
    \item $t_{.005,15}=2.947$
    \item $t_{.005,4}=4.604$
    \item $t_{.01,24}=2.492$
    \item $t_{.005,37}\approx2.712$
\end{enumerate}

\setcounter{enumi}{33}
\item
Given
\[
n=14,\qquad \bar{x}=8.48,\qquad s=.79,\qquad t_{.05,13}=1.771.
\]
\begin{enumerate}[label=\alph*.]
    \item A $95\%$ lower confidence bound is
    \[
    8.48-1.771\left(\frac{.79}{\sqrt{14}}\right)=8.48-.37=8.11.
    \]
    With $95\%$ confidence, the value of the true mean proportional limit stress of all such joints is greater than $8.11$ MPa. The sample observations must be from a normally distributed population.

    \item A $95\%$ lower prediction bound is
    \[
    8.48-1.771(.79)\sqrt{1+\frac{1}{14}}=8.48-1.45=7.03.
    \]
    If this bound is calculated for sample after sample, in the long run $95\%$ of these bounds will provide a lower bound for the corresponding future values of the proportional limit stress of a single joint of this type.
\end{enumerate}

\setcounter{enumi}{35}
\item
Given
\[
n=26,\qquad \bar{x}=370.69,\qquad s=24.36,\qquad t_{.05,25}=1.708.
\]
\begin{enumerate}[label=\alph*.]
    \item A $95\%$ upper confidence bound is
    \[
    370.69+(1.708)\left(\frac{24.36}{\sqrt{26}}\right)=370.69+8.16=378.85.
    \]

    \item A $95\%$ upper prediction bound is
    \[
    370.69+1.708(24.36)\sqrt{1+\frac{1}{26}}
    =370.69+42.45=413.14.
    \]

    \item Following the prediction interval argument, find the variance of
    \[
    \bar{X}-\bar{X}_{new}.
    \]
    Since
    \[
    V(\bar{X}-\bar{X}_{new})=V(\bar{X})+V(\bar{X}_{new})
    =V(\bar{X})+V\left(\frac{1}{2}(X_{27}+X_{28})\right),
    \]
    we get
    \[
    V(\bar{X})+V\left(\frac{1}{2}X_{27}\right)+V\left(\frac{1}{2}X_{28}\right)
    =\frac{\sigma^2}{n}+\frac{1}{4}\sigma^2+\frac{1}{4}\sigma^2
    =\sigma^2\left(\frac{1}{2}+\frac{1}{n}\right).
    \]
    Thus
    \[
    T=\frac{\bar{X}-\bar{X}_{new}}{s\sqrt{\frac{1}{2}+\frac{1}{n}}}
    \]
    has a $t$ distribution with $n-1$ df, so the new prediction interval is
    \[
    \bar{x}\pm t_{\alpha/2,n-1}s\sqrt{\frac{1}{2}+\frac{1}{n}}.
    \]
    For this situation,
    \[
    370.69\pm1.708(24.36)\sqrt{\frac{1}{2}+\frac{1}{26}}
    =370.69\pm30.53=(340.16,401.22).
    \]
\end{enumerate}

\setcounter{enumi}{41}
\item
\begin{enumerate}[label=\alph*.]
    \item $\chi^2_{.1,15}=22.307$ \quad (.1 column, 15 df row)
    \item $\chi^2_{.1,25}=34.381$
    \item $\chi^2_{.01,25}=44.313$
    \item $\chi^2_{.005,25}=46.925$
    \item $\chi^2_{.99,25}=11.523$ \quad (.99 column, 25 df row)
    \item $\chi^2_{.995,25}=10.519$
\end{enumerate}

\setcounter{enumi}{43}
\item
Here $n-1=8$, $\chi^2_{.025,8}=17.534$, and $\chi^2_{.975,8}=2.180$. Thus the $95\%$ confidence interval for $\sigma^2$ is
\[
\left(\frac{8(7.90)}{17.534},\frac{8(7.90)}{2.180}\right)=(3.60,28.98).
\]
The $95\%$ confidence interval for $\sigma$ is
\[
(\sqrt{3.60},\sqrt{28.98})=(1.90,5.38).
\]

\setcounter{enumi}{45}
\item
\begin{enumerate}[label=\alph*.]
    \item Using a normal probability plot, it is plausible that this sample came from a normal population distribution.
    \item With $s=1.579$, $n=15$, and $\chi^2_{.95,14}=6.571$, the $95\%$ upper confidence bound for $\sigma$ is
    \[
    \sqrt{\frac{14(1.579)^2}{6.571}}=2.305.
    \]
\end{enumerate}

\end{enumerate}

\section*{Chapter 8}

\begin{enumerate}[label=\textbf{\arabic*.}, leftmargin=*]

\setcounter{enumi}{3}
\item
Reject $H_0$ if and only if $P$-value $\le \alpha$.
\begin{enumerate}[label=\alph*.]
    \item Do not reject $H_0$, since $.084>.05$.
    \item Do not reject $H_0$, since $.003>.001$.
    \item Do not reject $H_0$, since $.498>.05$.
    \item Reject $H_0$, since $.084\le.10$.
    \item Do not reject $H_0$, since $.039>.01$.
    \item Do not reject $H_0$, since $.218>.10$.
\end{enumerate}

\setcounter{enumi}{7}
\item
\[
H_0:\mu=40\quad\text{versus}\quad H_a:\mu\ne40,
\]
where $\mu$ is the true average burn-out amperage for this type of fuse. The alternative reflects that departure from $\mu=40$ in either direction is a concern. A type I error would say that one of the two concerns exists, either $\mu<40$ or $\mu>40$, when the fuses are actually compliant. A type II error would be to fail to detect either concern when one exists.

\setcounter{enumi}{9}
\item
Let $\mu_1$ be the average amount of warpage for the regular laminate and $\mu_2$ the analogous value for the special laminate. Then
\[
H_0:\mu_1=\mu_2\quad\text{versus}\quad H_a:\mu_1>\mu_2.
\]
Type I error: conclude that the special laminate produces less warpage than the regular laminate when it actually does not. Type II error: conclude that there is no difference in the two laminates when, in reality, the special one produces less warpage.

\setcounter{enumi}{11}
\item
\begin{enumerate}[label=\alph*.]
    \item
    \[
    H_0:\mu=1300\quad\text{versus}\quad H_a:\mu>1300.
    \]

    \item When $H_0$ is true, $\bar{X}$ is normally distributed with mean $E(\bar{X})=\mu=1300$ and standard deviation
    \[
    \frac{\sigma}{\sqrt{n}}=\frac{60}{\sqrt{10}}=18.97.
    \]
    Values more contradictory to $H_0$ than $\bar{x}=1340$ would be anything above $1340$. Thus
    \[
    P\text{-value}=P(\bar{X}\ge1340\mid H_0\text{ true})
    =P\left(Z\ge\frac{1340-1300}{18.97}\right)
    =1-\Phi(2.11)=.0174.
    \]
    Since $.0174>.01$, $H_0$ would not be rejected at the $\alpha=.01$ significance level.

    \item When $\mu=1350$, $\bar{X}$ has a normal distribution with mean $1350$ and standard deviation $18.97$. To determine $\beta(1350)$, first find the threshold between $P$-value $\le\alpha$ and $P$-value $>\alpha$ in terms of $\bar{x}$:
    \[
    .01=P(\text{reject }H_0\mid H_0\text{ is true})=P(\bar{X}\ge\bar{x}\mid H_0\text{ is true})
    =1-\Phi\left(\frac{\bar{x}-1300}{18.97}\right).
    \]
    Hence
    \[
    \Phi\left(\frac{\bar{x}-1300}{18.97}\right)=.99
    \Rightarrow
    \frac{\bar{x}-1300}{18.97}=2.33
    \Rightarrow
    \bar{x}=1344.21.
    \]
    Thus we reject $H_0$ at the $\alpha=.01$ level if the observed value of $\bar{X}$ is at least $1344.21$. Finally,
    \[
    \beta(1350)=P(\text{do not reject }H_0\mid\mu=1350)
    =P(\bar{X}<1344.21\mid\mu=1350)
    \]
    \[
    =P\left(Z<\frac{1344.21-1350}{18.97}\right)
    \approx \Phi(-.31)=.3783.
    \]
\end{enumerate}

\setcounter{enumi}{29}
\item
The hypotheses are
\[
H_0:\mu=7.0\quad\text{versus}\quad H_a:\mu<7.0.
\]
In each case, the one-tail area to the left of the observed test statistic is needed.
\begin{enumerate}[label=\alph*.]
    \item $n=6\Rightarrow df=5$. From Table A.8,
    \[
    P(T\le-2.3\mid T\sim t_5)=P(T\ge2.3\mid T\sim t_5)=.035.
    \]
    Since $.035\le.05$, reject $H_0$ at the $\alpha=.05$ level.

    \item Similarly,
    \[
    P\text{-value}=P(T\ge3.1\mid T\sim t_{14})=.004.
    \]
    Since $.004<.01$, reject $H_0$.

    \item Similarly,
    \[
    P\text{-value}=P(T\ge1.3\mid T\sim t_{11})=.110.
    \]
    Since $.110\ge.05$, do not reject $H_0$.

    \item Here
    \[
    P\text{-value}=P(T\le .7\mid T\sim t_5),
    \]
    because this is a lower-tailed test. This is
    \[
    1-P(T>.7\mid T\sim t_5)=1-.258=.742.
    \]
    Since $.742>.05$, do not reject $H_0$.

    \item The observed test statistic is
    \[
    t=\frac{\bar{x}-\mu_0}{s/\sqrt{n}}=\frac{6.68-7.0}{.0820}=-3.90.
    \]
    Therefore,
    \[
    P\text{-value}=P(T\ge3.90\mid T\sim t_5)=.006.
    \]
    We would reject $H_0$ for any significance level at or above $.006$.
\end{enumerate}

\setcounter{enumi}{33}
\item
The data provided has
\[
n=6,\qquad \bar{x}=31.233,\qquad s=0.689.
\]
\begin{enumerate}[label=\alph*.]
    \item The hypotheses are
    \[
    H_0:\mu=30\quad\text{versus}\quad H_a:\mu>30.
    \]
    The one-sample $t$ statistic is
    \[
    t=\frac{31.233-30}{0.689/\sqrt{6}}=4.38.
    \]
    At $df=5$, the $P$-value is $P(T\ge4.38)=.004$. Since $.004\le.01$, reject $H_0$ at the $\alpha=.01$ level and conclude that the true average stopping distance exceeds $30$ ft.

    \item Using statistical software, the power of the test with $n=6$, $\alpha=.01$, and $\sigma=.65$ is $.660$ when $\mu=31$ and $.998$ when $\mu=32$. The corresponding type II error probabilities are
    \[
    \beta(31)=1-.660=.340,
    \qquad
    \beta(32)=1-.998=.002.
    \]

    \item Changing $\sigma$ from $.65$ to $.80$ gives
    \[
    \beta(31)=1-.470=.530,
    \qquad
    \beta(32)=1-.975=.025.
    \]
    Type II error is still less likely when the true value of $\mu$ is farther from $\mu_0$. With the larger standard deviation, random sampling variation makes it more likely for the sample mean to deviate substantially from $\mu$, and in particular tricks us into thinking $H_0$ might be plausible even though it is false. Thus $\beta$ increases from $.340$ to $.530$ when $\mu=31$ and from $.002$ to $.025$ when $\mu=32$.

    \item According to statistical software, $n=9$ will suffice to achieve $\beta=.1$ (power $=.9$) when $\alpha=.01$, $\mu=31$, and $\sigma=.65$.
\end{enumerate}

\setcounter{enumi}{37}
\item
$\mu$ is the true average percentage of organic matter in this type of soil, and the hypotheses are
\[
H_0:\mu=3\quad\text{versus}\quad H_a:\mu\ne3.
\]
With $n=30$ and assuming normality, use the $t$ test:
\[
t=\frac{\bar{x}-3}{s/\sqrt{n}}
=\frac{2.481-3}{.295}=
\frac{-.519}{.295}=-1.759.
\]
At $df=29$,
\[
P\text{-value}=2P(T>1.759)=2(.041)=.082.
\]
At significance level $.10$, since $.082\le.10$, reject $H_0$ and conclude that the true average percentage of organic matter is something other than $3$. At significance level $.05$, do not reject $H_0$.

\setcounter{enumi}{39}
\item
The hypotheses are
\[
H_0:\mu=48\quad\text{versus}\quad H_a:\mu>48.
\]
Using the one-sample $t$ procedure,
\[
t=\frac{51.3-48}{1.2/\sqrt{10}}=8.7.
\]
With $df=9$,
\[
P(T\ge8.7)\approx0.
\]
Hence reject $H_0$ at any reasonable significance level and conclude that the true average strength for the WSF/cellulose composite definitely exceeds $48$ MPa.

\end{enumerate}

\end{document}
